% !TeX root = main.tex
\chapter{Theoretical Background}

\section{Machine Translation}

\subsection{History}

\quad Machine translation is one of the oldest goals of natural language
processing, and its history is easiest to follow as a sequence of paradigms.
Each paradigm gave a different answer to the same question: where should the
knowledge needed for translation come from? Rule-based systems encoded it by
hand, statistical systems estimated it from parallel data, and neural systems
learn it end to end inside a single trainable model.

Rule-based machine translation dominated the early decades of the field.
Systems from this period combined bilingual dictionaries, hand-written
grammatical rules, and transfer procedures designed by linguists. The
best-known representative is SYSTRAN, which was used for many years in real
institutional and commercial settings. The strength of the rule-based approach
was control: because every rule was written by a person, the behavior of the
system could be inspected, explained, and corrected. Its weakness was cost:
every new language pair and every new domain required additional expert work.

Statistical machine translation replaced hand-written rules with probabilities
estimated from bilingual corpora: researchers collected aligned sentence pairs
and let the system learn word and phrase correspondences from the data. The
IBM alignment models
introduced the mathematical core of this approach, and phrase-based systems
later turned it into the dominant technology of the 2000s. Statistical systems
clearly outperformed rule-based ones whenever enough parallel data existed,
but they still translated through a pipeline of separately trained components,
which limited how much context each decision could use.

Neural machine translation changed this picture by learning the entire mapping
from the source sequence to the target sequence in one end-to-end model. In
the early neural period, the standard architecture was an encoder-decoder
built from recurrent networks. The system that best marked the maturity of
this paradigm was Google's Neural Machine Translation system, GNMT
\cite{wu2016gnmt}, which showed at production scale that neural models could
overtake strong phrase-based baselines.

The current period began with the Transformer architecture
\cite{vaswani2017transformer}. By replacing recurrence with self-attention,
the Transformer trains efficiently in parallel and models long-distance
dependencies directly, because any token can attend to any other token in the
sequence. These properties quickly made it the dominant architecture in
machine translation.

\subsection{SoTA}

\quad In 2026, the state of the art in machine translation is no longer tied
to a single model name. Different systems lead in different settings, and the
strongest results come from a family of very large multilingual
Transformer-based models, often combined with LLM-style adaptation or
prompting strategies. A good reference point for this direction is the No
Language Left Behind project and its NLLB-200 model, which showed that one
multilingual system can cover around two hundred languages while keeping
strong quality, including for many low-resource pairs \cite{team2024nllb}.

Recent shared-task results show that the field has entered a more mixed phase.
The WMT24 General Machine Translation task makes two points at the same time:
large language models are now central to machine translation, and the task is
still not solved \cite{wmt24general}. Very strong results are possible, but
they depend on the language pair, the domain, the model size, the adaptation
strategy, and the evaluation protocol. No single system wins in every setting.

The way quality is judged has also become broader: beyond automatic scores on
one benchmark, researchers look at multilingual robustness, domain transfer,
stability on low-resource languages, and the gap between automatic metrics and
real usefulness. Modern machine translation research is therefore less about
finding one new architecture and more about combining strong pretrained
models, multilingual data, and careful tuning and evaluation.

\subsection{Code translation}

\quad Code translation applies the machine translation idea to programming
languages, but the two tasks differ in one essential point: the standard of
correctness. A natural language sentence usually has several acceptable
translations, and small variations in word choice rarely make the output
wrong. A program is judged much more strictly. A translated program that does
not compile, crashes at runtime, or computes a different result is not an
imperfect translation but a wrong one, no matter how similar it looks to the
reference.

The structure of the data is also different. Natural language tolerates
ambiguity, redundancy, and stylistic variation. Programming languages are
formal: a translation must preserve the semantics of the source program while
respecting the syntax, the scoping rules, the type system, the standard
library, and the conventions of the target language. A single changed
identifier or a shifted loop bound can silently change the behavior of an
otherwise plausible program. Code translation is therefore not only sequence
generation; the model must also track how values flow through the program.

The research history of code translation is shorter than that of general
machine translation. Early systems were closer to transpilers and compiler
engineering than to learned models: they worked well when the mapping between
the two languages could be described by rules, but scaled poorly once the
correspondence became less direct.

The modern phase started when pretrained neural models for source code became
available. An important milestone was TransCoder, which showed that
translation between programming languages is possible even without parallel
programs, using unsupervised training \cite{roziere2020transcoder}. Pretrained
encoder-decoder models such as PLBART then showed that code generation and
translation benefit from pretraining on large amounts of source code and
technical natural language \cite{ahmad2021plbart}. The overall direction
mirrors what happened in general machine translation, with one important
difference: in code translation, evaluation must go beyond text similarity,
because a translation is useful only if it preserves the behavior of the
original program.

\section{AVATAR Dataset}

\quad AVATAR is the central dataset of this thesis. It was built specifically
for program translation between Java and Python \cite{ahmad2023avatar}. Unlike
the large code corpora used for pretraining, AVATAR is a parallel corpus: for
each programming problem it contains a Java solution and a Python solution
that solve the same task. This alignment provides the direct source-target
relation needed for supervised translation experiments and for a controlled
comparison between models.

The examples come from competitive-programming platforms and similar sources,
so each program is short, self-contained, and focused on a single algorithmic
idea. This format suits translation research well: each pair shows how the
same algorithm is expressed in the idioms of the two languages, without the
project-level dependencies of full software repositories. The dataset comes
with a fixed train, validation, and test split, and all models in this thesis
are trained and evaluated on the same split, so the reported differences can
be attributed to the models rather than to the data.

The property that matters most for this thesis is that AVATAR connects each
problem to input-output test cases. Text similarity alone cannot establish
whether a translation is correct: a generated program can be close to the
reference at the token level and still fail to compile or print a wrong
answer. The test cases make it possible to compile and execute every generated
translation and to check its actual behavior, and they are the basis of the
execution-based metrics used in this work.

AVATAR also has clear limitations. It covers a single language pair, and its
programs reflect competitive-programming style rather than production
software, so its results do not automatically transfer to every code
translation scenario. Within these limits, it remains one of the most useful
resources for controlled Java-Python translation experiments.

\section{Evaluation Metrics}

\quad The evaluation in this thesis uses two families of metrics. The first
family measures how similar the generated program is to the reference
translation, as text or as structure: BLEU, Syntax Match, Dataflow Match, and
CodeBLEU. The second family ignores the reference text and asks what the
generated program actually does when it is compiled and executed: Execution
Accuracy and Computational Accuracy. The two families answer different
questions, and the experiments in this thesis show that they can disagree
strongly: a system can produce code that looks very close to the reference and
still fails the tests. For this reason, every model is reported on both
families.

\subsection{BLEU}

\quad BLEU is the classical metric of machine translation
\cite{papineni2002bleu}. It compares the candidate translation with the
reference at the level of n-grams, which are sequences of $n$ consecutive
tokens. For each order $n$ from 1 to 4, BLEU computes a modified precision
$p_n$: the fraction of candidate n-grams that also appear in the reference,
where each reference n-gram can be matched at most as many times as it occurs
in the reference. The four precisions are combined by a geometric mean, and
the result is multiplied by a brevity penalty $\mathrm{BP}$ that lowers the
score when the candidate is shorter than the reference:
\[
\text{BLEU} = \mathrm{BP} \cdot \exp\left(\sum_{n=1}^{4} \tfrac{1}{4} \log p_n\right),
\qquad
\mathrm{BP} = \min\left(1,\; e^{\,1 - r/c}\right),
\]
where $r$ is the total reference length and $c$ is the total candidate length.
In the convention used here, the score is scaled to the range from 0 to 100
and computed over the whole test corpus rather than per example.

For code, BLEU is applied to the token-spaced representation of the programs,
so it measures how many token patterns the translation shares with the
reference solution. It is a useful metric, but it does not reveal the whole information about translation. 
However, BLEU has no notion
of program structure or meaning. Renaming one variable lowers the score
without changing the behavior, while replacing \texttt{<} with \texttt{<=},
which can break the program, costs almost nothing.

\subsection{Syntax Match}

\quad Syntax Match (SM) is the structural component of CodeBLEU
\cite{ren2020codebleu}. Both the candidate and the reference program are
parsed into abstract syntax trees, the trees are decomposed into subtrees, and
the score is the fraction of reference subtrees that also occur in the
candidate tree. Intuitively, it asks whether the translation is built from the
same syntactic constructions as the reference: the same loop nesting, the same
conditional shapes, the same expression patterns.

The motivation is that source code is hierarchical, not linear. Two programs
can share most of their tokens and still have different structure, and the
structure carries meaning: moving a statement into or out of a loop body
changes the behavior completely while changing very little at the token level.
Syntax Match captures this kind of agreement better than n-gram overlap. Its
limitation is that it still compares against a single reference, so a correct
program written with a different structure, for example a \texttt{while} loop
instead of a \texttt{for} loop, is penalized even when its behavior is
identical.

\subsection{Dataflow Match}

\quad Dataflow Match (DM) is the semantic component of CodeBLEU
\cite{ren2020codebleu}. From each program, a data-flow graph is extracted: the
nodes are variables and the edges record where each value comes from, for
example which assignment defines the variable that a later expression reads.
The score is the fraction of reference data-flow edges that are matched in the
candidate program.

This metric looks one level deeper than syntax. A translation can use the
right constructions in the right order and still connect the wrong values, for
instance by reading a loop counter where the accumulated sum was intended.
Such a mistake barely affects BLEU or Syntax Match, but it breaks the
dependency structure that Dataflow Match checks. Because variable names are
abstracted away during matching, the metric is also robust to renaming, which
plain token overlap is not. Among the text-based metrics it is the closest to
semantics, although it still cannot confirm that the program computes the
correct result.

\subsection{CodeBLEU}

\quad CodeBLEU combines the previous signals into one score designed for code
\cite{ren2020codebleu}. It is a weighted sum of four components:
\[
\text{CodeBLEU} = \alpha \cdot \text{BLEU} + \beta \cdot \text{BLEU}_{w} +
\gamma \cdot \text{SM} + \delta \cdot \text{DM},
\]
where $\text{BLEU}_{w}$ is a weighted n-gram match that gives the keywords of
the language more importance than ordinary tokens. In the standard
configuration used in this thesis, the four weights are equal,
$\alpha = \beta = \gamma = \delta = 0.25$.

The intent is to judge code the way a reviewer would: partly by its tokens,
partly by its structure, and partly by how data moves through it. CodeBLEU is
usually more informative than plain BLEU for code tasks. The averaging also
has a side effect that matters when reading result tables: a weak component
can be hidden by strong ones, so the individual components are reported
alongside the combined value.

\subsection{Execution Accuracy}

\quad Execution Accuracy (EA) is the first execution-based metric
\cite{ahmad2023avatar}. Every generated translation is converted back into
normal source code and actually run on the test inputs of its problem. A
program counts as executable when it passes the language processing stage,
which is compilation for Java and a syntax check for Python, and then runs on
the test inputs without raising a runtime error. Execution Accuracy is the
percentage of test problems whose translation is executable in this sense.
Programs that run but print a wrong answer still count, because this metric
measures only the ability to run.

Execution Accuracy is a much stricter filter than any similarity score,
because it is applied by the compiler and the runtime rather than by
comparison with a reference. It catches missing declarations, type errors,
syntax mistakes, and crashes, which are exactly the failures that text metrics
tend to overlook. Its limit is that it says nothing about what the program
computes, so a high Execution Accuracy alone does not mean that the
translations are correct.

\subsection{Computational Accuracy}

\quad Computational Accuracy (CA) is the strictest metric in this thesis and
its main criterion for comparing models. A translation is counted as correct
only when it compiles, runs on every available test case of its problem, and
produces the expected output on all of them; numeric outputs are compared with
a small tolerance. The metric is the percentage of test problems for which
this happens. The same idea was introduced for code translation by TransCoder
\cite{roziere2020transcoder} and is used by the AVATAR benchmark
\cite{ahmad2023avatar}.

Computational Accuracy is the closest automatic proxy to the real goal of the
task: a translated program that behaves exactly like the source program. It
subsumes the other metrics in one direction only: a program that passes all
tests usually also has reasonable similarity scores, but high BLEU or CodeBLEU
values do not imply high Computational Accuracy. This gap is a recurring
observation in the results of this thesis, and it is the main reason why
execution-based evaluation is treated as mandatory for code translation.
